\chapter{Schemas and Reproducibility}

\section{Common Data Record}\label{app:data-record}

The serialized interchange object contains the following fields. Optional empty
fields may be omitted.

\begin{longtable}{p{.24\textwidth}p{.18\textwidth}p{.47\textwidth}}
\caption{Common data-record schema.}\\
\toprule
\textbf{Field} & \textbf{Type} & \textbf{Meaning}\\
\midrule
\endfirsthead
\toprule
\textbf{Field} & \textbf{Type} & \textbf{Meaning}\\
\midrule
\endhead
\texttt{\_type} & string & \texttt{data}, \texttt{session\_start},
\texttt{session\_end}, or \texttt{error}.\\
\texttt{session\_id} & string & Identifier shared by one runtime session.\\
\texttt{record\_id} & string & Stable observation identifier within the
session.\\
\texttt{source\_type} & string & \texttt{topic}, \texttt{service}, or
\texttt{action}.\\
\texttt{source\_name} & string & Fully qualified ROS~2 resource name.\\
\texttt{source\_node} & string array & Discovered publisher or server node
names, when available.\\
\texttt{msg\_type} & string & ROS~2 interface type.\\
\texttt{phase} & string/null & Service \texttt{request}/\texttt{response} or
action \texttt{feedback}/\texttt{status}; absent for normal topics.\\
\texttt{timestamp} & number & Collector wall-clock receive time in seconds.\\
\texttt{ros\_timestamp} & object/null & Optional \texttt{sec}/\texttt{nanosec}
timestamp from the observation.\\
\texttt{data} & object & Recursively JSON-compatible payload.\\
\texttt{metadata} & object & Per-source sequence and processing context.\\
\bottomrule
\end{longtable}

The following is a complete topic-record example:

\begin{lstlisting}[language={},caption={Example common data record.}]
{
  "_type": "data",
  "session_id": "20260727_101530_a1b2c3d4",
  "record_id":
    "20260727_101530_a1b2c3d4:topic:/cmd_vel:-:18",
  "source_type": "topic",
  "source_name": "/cmd_vel",
  "source_node": ["/controller_server"],
  "msg_type": "geometry_msgs/msg/Twist",
  "timestamp": 1785147330.274,
  "data": {
    "linear": {"x": 0.22, "y": 0.0, "z": 0.0},
    "angular": {"x": 0.0, "y": 0.0, "z": 0.14}
  },
  "metadata": {
    "seq": 18,
    "transformers_applied": ["FieldExtractor"]
  }
}
\end{lstlisting}

\section{Deployment Specification Example}\label{app:deployment-spec}

This complete example places ROS~2 collection on \texttt{robot} and speed
checking on \texttt{verifier}. The generator creates
\texttt{robot.yaml} for \texttt{monitor\_node} and
\texttt{verifier.yaml} for \texttt{node\_runner};
Listings~\ref{lst:gen-robot} and~\ref{lst:gen-verifier} show both files
exactly as generated from this specification.

\begin{lstlisting}[language={},caption={Two-host deployment specification.}]
{
  "hosts": [
    {
      "id": "robot",
      "runtimes": [
        {
          "id": "ros2_robot",
          "kind": "ros2",
          "sources": [
            {
              "id": "cmd",
              "source_kind": "topic",
              "name": "/cmd_vel",
              "interface": "geometry_msgs/msg/Twist"
            }
          ]
        },
        {
          "id": "monitor",
          "kind": "monitor",
          "subscribe": ["cmd"]
        }
      ]
    },
    {
      "id": "verifier",
      "runtimes": [
        {
          "id": "speed_converter",
          "kind": "converter",
          "class_path": "custom.speed:CmdVelSpeedConverter",
          "input_from": ["cmd"],
          "output_to": "speed_dsl"
        },
        {
          "id": "speed_property",
          "kind": "verdict_service",
          "class_path": "custom.threshold:ThresholdVerdict",
          "input_from": ["speed_dsl"],
          "output_to": [
            {"kind": "stdout"},
            {"kind": "file",
             "path": "verdicts_{session_id}.jsonl"}
          ],
          "property_id": "cmd_speed_limit",
          "field": "speed",
          "op": "<=",
          "threshold": 0.5
        }
      ]
    }
  ],
  "links": [
    {
      "id": "robot_records",
      "from_host": "robot",
      "to_host": "verifier",
      "from_runtime": "monitor",
      "to_runtime": "speed_converter",
      "payload": "records",
      "transport": {
        "kind": "mqtt",
        "broker": "192.0.2.10",
        "port": 1883,
        "qos": 1,
        "topic_prefix": "monitor/"
      }
    }
  ]
}
\end{lstlisting}

\begin{lstlisting}[language={},caption={Generated \texttt{robot.yaml},
started with \texttt{monitor\_node}.},label={lst:gen-robot}]
monitor:
  output_dir: ./output/robot
  session_id_prefix: robot
topics:
- name: /cmd_vel
  type: geometry_msgs/msg/Twist
  exporters:
  - type: mqtt
    broker: 192.0.2.10
    port: 1883
    topic_prefix: monitor/
    qos: 1
\end{lstlisting}

\begin{lstlisting}[language={},caption={Generated \texttt{verifier.yaml},
started with \texttt{node\_runner}.},label={lst:gen-verifier}]
monitor:
  output_dir: ./output/verifier
inputs:
- id: robot_records
  payload: records
  type: mqtt
  broker: 192.0.2.10
  port: 1883
  topic_filter: monitor/#
  qos: 1
converters:
- id: speed_converter
  type: custom.speed:CmdVelSpeedConverter
verdict_services:
- id: speed_property
  type: custom.threshold:ThresholdVerdict
  params:
    property_id: cmd_speed_limit
    field: speed
    op: <=
    threshold: 0.5
  exporters:
  - type: stdout
  - type: file
    path: verdicts_{session_id}.jsonl
links:
- from: source:/cmd_vel
  to: converter:speed_converter
- from: converter:speed_converter
  to: verdict:speed_property
\end{lstlisting}

The single MQTT link in the specification produces both endpoints: the
publishing prefix \texttt{monitor/} in \texttt{robot.yaml} and the
subscription filter \texttt{monitor/\#} in \texttt{verifier.yaml}. The
converter and verdict-service settings of the \texttt{verifier} host become
the \texttt{converters}, \texttt{verdict\_services}, and \texttt{links}
sections of \texttt{verifier.yaml}.

\section{Plug-in Contracts}\label{app:plugin-contracts}

A converter subclasses \texttt{DataConverter} and implements:

\begin{lstlisting}[language=Python]
def convert(self, record: DataRecord) -> Any | None:
    ...
\end{lstlisting}

A verdict service subclasses \texttt{VerdictService} and implements:

\begin{lstlisting}[language=Python]
def evaluate(self, dsl_record: Any) -> Verdict | None:
    ...
\end{lstlisting}

A custom transformer subclasses \texttt{Transformer} and implements:

\begin{lstlisting}[language=Python]
def transform(self, record: DataRecord) -> DataRecord | None:
    ...
\end{lstlisting}

Returning \texttt{None} drops the record. Built-in transformers are selected
by short name; a custom transformer uses the same
\texttt{module.path:ClassName} import path as converters and verdict
services.

The converter must preserve the input record identifier in its DSL record when
the resulting verdict needs attribution. The verdict service returns the common
\texttt{Verdict} object; output destinations are configured separately.

\section{Repository}\label{app:repository}

The implementation, tests, documentation, and case-study material are available
at:

\begin{center}
\url{https://github.com/Jin-Shikai/ROS2-Monitor-Infra}
\end{center}
